\section{Discussion: what to log and what to watch}
\label{sec:discussion}

\paragraph{Action rules detect successful attacks; refusal logging detects attempts.} Against models
that act on injections, rules over the agent's actions and output caught all 269 successful attacks
in our ten scenarios. Against frontier models the same rules have almost nothing to see, because the
attack fails before any action. The attempt still leaves evidence, in a place most applications do
not log: a provider refusal (an empty answer with a stop reason) or a refusal in the model's own words.
A SOC that wants to know it is being targeted, and not only that it has been breached, needs the
application to log provider refusals as events and to keep the answer text for a rule to read. In
our runs this raised coverage of frontier attack sessions from 13\% to 70\% at an 8\% benign rate.

\paragraph{Attempt, compromise and harm are different detections.} Input rules see attempts and are
easily evaded; rules over actions and output see compromise and harm, but only when the model acts;
the refusal signal sees attempts that failed. A single ``coverage'' number mixes these, so we report
them separately, and a SOC should route them separately, treating attempts as threat intelligence and
compromises as incidents.

\paragraph{Log what the model saw.} In our experiments an application preview shorter than the model's
input hid a followed instruction from the log, a guardrail blocked nothing past the first 1{,}000
characters of a user message and nothing in tool results, and the classifiers did not recognise the
N1 payload at any offset. The cheapest fix is in the application. It should log the full tool result,
or its length together with the head, the tail and a hash of the rest, so that a rule can at least
tell that something was not inspected. A guardrail that raises no alert may not have inspected the
whole input.

\paragraph{Watch the answer, not only the actions.} In the nine models that relayed, declining to act
on an injection usually meant passing the attacker's request on to the user as a legitimate next
step, in 61\% to 100\% of runs (Section~\ref{sec:relay-models}). The attack then depends on the human,
and the only telemetry is the answer text. \llmd{026} is a first, narrow rule for this: an internal
resource name and an outside address in the same answer, without a warning. A production version
needs the organisation's own list of internal resource names, which most SOCs already keep for data
loss prevention.

\paragraph{Content signatures as triage context.} The input rules rarely fire on benign public
traffic and miss reworded attacks; the classifiers recall far more direct injections but little
indirect injection in documents. Both help with triage and with explaining an alert, and neither
should be the only line of detection. Their alerts should be routed at a lower severity than
behavioural and refusal signals.
